\documentclass[a4paper, titlepage]{book}

\usepackage{amsmath}
\numberwithin{equation}{section}
\usepackage{amssymb}
\usepackage{csquotes}
\MakeOuterQuote{"}
\usepackage[useregional]{datetime2}
\DTMsetstyle{iso}
\usepackage{hyperref}
\hypersetup{
    colorlinks = true,
    allcolors = blue
}
\usepackage{indentfirst}
\usepackage[skip=5pt, indent=30pt]{parskip}
\usepackage{pgf}
\usepackage{physics2}
\usepackage{siunitx}
\sisetup{
    input-digits = 0123456789\pi,
    inter-unit-product = \cdot,
}


\title{
    Eric Zhang's Physics Notes \\
    \large Volume 3: Waves and Oscillations
}
\author{Eric Zhang}
\date{\today}


\begin{document}
\frontmatter
\maketitle
\tableofcontents


\mainmatter

\chapter{Springs and Pendulums}

\section{Pendulums}

Simple pendulum have a massles rod/string;
\(\therefore\) they have no moment of inertia.

By Newton's 2nd Law:
\begin{align}
    F = ma \\
    -mg \sin \theta = m \ell \ddot{\theta}
\end{align}

For angles less than approximately \qty{14}{\degree},
the small angle approximation results in a error of less than 1\% in the gravitational force:
\begin{align}
    \sin\theta \approx \theta \\
    \therefore -mg\theta            &\approx m \ell \ddot{\theta} \\
    \ddot{\theta} + \omega^2 \theta &\approx 0,~\omega = \sqrt{\frac{g}{\ell}}
\end{align}

For physical pendulums, rotational inertia \(\ne 0\).
\begin{align}
    -mgd \sin\theta                &=       I \ddot{\theta} \\
    \ddot{\theta} + \omega^2\theta &\approx 0,~\omega = \sqrt{\frac{mgd}{I}}
\end{align}
where \(d\) is the distance between center of mass and pivot.

\section{Simple Springs}
Simple springs operate are governed by only Newton's 2nd Law and Hooke's Law.

\begin{align}
    F              &= ma = m\ddot{x} \\
    F              &= -kx \\
    m\ddot{x}      &= -kx \\
    m\ddot{x} + kx &= 0
\end{align}

Solution forms:
\begin{itemize}
    \item \(x(t) = A\cos\left(\omega_0 t\right) + B\sin\left(\omega_0 t\right)\)
    \item \(x(t) = A\cos\left(\omega_0 t + \phi\right)\)
    \item \(x(t) = z_0 e^{i \omega_0 t} + z^*_0 e^{i \omega_0 t}\)
    \item \(\tilde{x}(t) = A e^{i \phi} e^{i \omega_0 t}\),
          where \(\Re[\tilde{x}(t)] = x(t)\)
\end{itemize}

\begin{equation}
    \tilde{v}(t) = \frac{d}{dt} \tilde{x}(t)
    = i\omega_0 A e^{i\phi} e^{\omega_0 t}
    = i\omega_0 \tilde{x}(t)
\end{equation}
The factor of \(i\) gives a \qty{90}{\degree} phase shift.

\begin{equation}
    \tilde{a}(t) = \frac{d}{dt} \tilde{v}(t)
    = -\omega_0^2 A e^{i\phi} e^{\omega_0 t}
    = -\omega_0^2 \tilde{x}(t)
\end{equation}
The minus signifies a \qty{180}{\degree} phase shift.


\chapter{Electrical Circuits}


\chapter{Systems of Oscillators}


\chapter{Transverse Waves}


\chapter{Longitudinal Waves}


\chapter{Dispersion}


\chapter{Electromagnetic Waves}


\appendix

\chapter{Complex and Imaginary Numbers} \label{section:complex-and-imaginary-numbers}

\(i = \sqrt{-1}\) is the \textbf{imaginary unit}.

\textbf{Imaginary numbers} are the set of all numbers that can be written in the form \(bi\),
where \(b \in \mathbb{R}\).

\textbf{complex numbers}, denoted by \(\mathbb{C}\),
are the set of all numbers that can be written in the form \(a + bi\),
where \(a,~b \in \mathbb{R}\).

\(a\) is known as the \textbf{real part},
and \(b\) is known as the \textbf{imaginary part}.
Note that this means the imaginary part is \emph{real}.
Given the complex number \(z = a + bi\),
\(\mathrm{Re}[z] = \Re[z] = a\), and \(\mathrm{Im}[z] = \Im[z] = b\).

The \textbf{complex conjugate} of \(z = a + bi\) is given by \(z^* = a - bi\).

\section{Complex Plane}

Complex numbers can be represented as 2D vectors on a plane,
with the x-axis being the real part (i.e., \(a\)),
and the y-axis being the imaginary part (i.e., \(b\)).

Given a complex number \(z\),
its magnitude \(\left\|z\right\|\) is given by Pythagorean's Theorem:
\begin{equation}
    \left\|z\right\| = \sqrt{\Re[z]^2 + \Im[z]^2} = \sqrt{z z^*}
\end{equation}

Its angle \(\phi\) is derived through trigonometry
(be careful about the general limitations of \(\arctan\) on calculators):
\begin{equation}
    \phi = \arctan\left(\frac{\Im[z]}{\Re[z]}\right)
\end{equation}

\section{Euler's Formula}

Euler's formula:
\begin{equation}
    e^{i \theta} = \cos\theta + i \sin\theta
\end{equation}

\begin{equation}
    \therefore e^{a + bi} = e^a e^{bi} = e^a \left(\cos b + i \sin b\right)
\end{equation}

Oftentimes, working with the complex exponential is easier than the trigonometric functions.

\section{Complex-Valued Function}

A \textbf{complex-valued function} is a function that returns a complex number.
It is denoted with a overhead tilde, e.g., \(\tilde{f}\left(x\right)\).


\chapter{Fourier Analysis}

\section{Fourier Series}

\section{Fourier Transform}


\backmatter

\chapter{References}


\end{document}